%!TEX encoding = UTF8
%!TEX root = 0-notes.tex

\chapter{Dénombrement}
\label{chap:2}

Ou comment compter des objets.
Cardinal.

\section{Ensembles finis}

\subsection{Sous-ensembles et combinaisons}

\subsection{Coefficients binomiaux}

\subsection{Formule du crible}

\section{Applications entre ensembles finis}

\subsection{Injectivité, surjectivité, bijectivité}

\subsection{Stars and bars}

\section{Ensembles infinis}

\subsection{Comparaisons d'ensembles et dénombrabilité}

\subsection{Indénombrabilité et argument diagonal de Cantor}

% corollaire : X et P(X) n'ont pas le même cardinal

\section{Relations d'équivalences}

\subsection{Relation binaire}

\subsection{Ordre partiel, ordre total}

\section{Applications : probabilités}



\section*{Exercices supplémentaires}
\addcontentsline{toc}{section}{Exercices supplémentaires}

\exe{,difficulty=3}{
	exe2
}{exe:2}{
	sol2
}

%\section*{Solutions}
%\addcontentsline{toc}{section}{Solutions}
%
%\shipoutAnswer



